\documentclass[11pt]{article}

\usepackage[margin=1in]{geometry}
\usepackage{amsmath,amssymb,amsthm}
\usepackage{booktabs}
\usepackage{enumitem}
\usepackage{hyperref}
\usepackage{xcolor}

\hypersetup{colorlinks=true,linkcolor=blue,urlcolor=blue,citecolor=blue}
\emergencystretch=2em

\newtheorem{theorem}{Theorem}
\newcommand{\trnorm}[1]{\lVert#1\rVert_1}
\newcommand{\opnorm}[1]{\lVert#1\rVert_{\mathrm{op}}}

\title{A Two-by-Two Counterexample to\\Zou's Heinz--Heron Norm Conjecture}
\author{Math discovery packet}
\date{August 13, 2026}

\begin{document}
\maketitle

\begin{center}
\fcolorbox{black}{gray!10}{\parbox{0.91\textwidth}{
\textbf{Outcome.}  The conjecture in Remark 4.2 of L.~Zou,
\emph{J. Math. Inequal.} 7 (2013), 389--397, is false.  A real
\(2\times2\) counterexample uses the trace norm, \(\nu=17/20\), and
data whose entries and logarithmic spectra are all rational tenths.
An exact-rational interval certificate proves a strict gap greater than
\(1001/10^6\).}}
\end{center}

\section{The conjecture and a transcription issue}

For positive matrices \(A,B\), an arbitrary matrix \(X\), and
\(0\leq\nu\leq1\), set
\[
 \alpha(\nu)=1-4(\nu-\nu^2)=(2\nu-1)^2.
\]
Zou conjectured that every unitarily invariant norm satisfies
\begin{align}
 \frac12\left\lVert A^\nu XB^{1-\nu}+A^{1-\nu}XB^\nu\right\rVert
 &\leq (1-\alpha(\nu))\left\lVert A^{1/2}XB^{1/2}\right\rVert
 \notag\\[-2pt]
 &\quad+\alpha(\nu)\left\lVert\frac{AX+XB}{2}\right\rVert .
 \label{ZH}
\end{align}
This is inequality (4.2) and the Conjecture in Remark 4.2 of the
included 2013 primary paper.

The lane target is Y.~Kapil, C.~Conde, M.~S.~Moslehian, M.~Singh, and
M.~Sababheh, \emph{Norm inequalities related to the Heron and Heinz
means}, arXiv:1709.09813.  In its ``On Zou's questions'' paragraph it
prints the second Heinz term as \(B^\nu XA^{1-\nu}\).  Comparison with
Zou's original display shows that this is a transcription error.  The
crossed version is already false at \(\nu=1/2\) by elementary matrix-unit
examples.  The result below addresses the mathematically intended and
correctly sourced conjecture \eqref{ZH}, not that typo.

\section{Explicit counterexample}

Let \(\trnorm{M}=\operatorname{tr}((M^*M)^{1/2})\), a unitarily
invariant norm.  Put
\[
 \nu=\frac{17}{20},\qquad s=2\nu-1=\frac7{10},
 \qquad \alpha(\nu)=s^2=\frac{49}{100},
\]
and define
\[
 A=\begin{pmatrix}e^{3/5}&0\\0&e^{13/5}\end{pmatrix},\qquad
 B=\begin{pmatrix}1&0\\0&e^{9/5}\end{pmatrix},\qquad
 T=\frac1{10}\begin{pmatrix}5&-8\\-1&4\end{pmatrix}.
\]
Finally, let
\[
 X=A^{-1/2}TB^{-1/2}.
\]
The matrices \(A,B\) are positive definite and \(X\) is a real
\(2\times2\) matrix.

\begin{theorem}
For these data, inequality \eqref{ZH} fails for the trace norm.  In fact,
its left-hand side exceeds its right-hand side by more than
\(1001/10^6\).
\end{theorem}

\begin{proof}
Write
\[
 h=\left(\frac3{10},\frac{13}{10}\right),\qquad
 k=\left(0,\frac9{10}\right),\qquad d_{ij}=h_i-k_j.
\]
Since \(A_{ii}=e^{2h_i}\) and \(B_{jj}=e^{2k_j}\), direct entrywise
calculation gives
\begin{align*}
 Y&:=\frac12\left(A^\nu XB^{1-\nu}+A^{1-\nu}XB^\nu\right),
 &Y_{ij}&=\cosh(sd_{ij})T_{ij},\\
 A^{1/2}XB^{1/2}&=T,\\
 Z&:=\frac{AX+XB}{2},
 &Z_{ij}&=\cosh(d_{ij})T_{ij}.
\end{align*}
Thus, explicitly,
\[
Y=\begin{pmatrix}
 \frac12\cosh(21/100)&-\frac45\cosh(21/50)\\
 -\frac1{10}\cosh(91/100)&\frac25\cosh(7/25)
\end{pmatrix},
\]
and
\[
Z=\begin{pmatrix}
 \frac12\cosh(3/10)&-\frac45\cosh(3/5)\\
 -\frac1{10}\cosh(13/10)&\frac25\cosh(2/5)
\end{pmatrix}.
\]

For a real \(2\times2\) matrix \(M\), its singular values give the
identity
\begin{equation}
 \trnorm{M}^2=\lVert M\rVert_F^2+2|\det M|.
 \label{traceformula}
\end{equation}
The exact interval calculation described below proves
\begin{equation}
 \trnorm{Y}>1.178,\qquad
 \trnorm{T}<1.1402,\qquad
 \trnorm{Z}<1.2153.
 \label{bounds}
\end{equation}
It follows that the right-hand side of \eqref{ZH} is strictly less than
\[
 \frac{51}{100}(1.1402)+\frac{49}{100}(1.2153)
 =1.176999,
\]
whereas its left-hand side is greater than \(1.178\).  Their difference
is therefore greater than
\(1.178-1.176999=0.001001=1001/10^6\).
\end{proof}

\section{Exact certificate for the three norm bounds}

For completeness, here is a reproducible proof of \eqref{bounds}.  If
\(x\in\mathbb Q\) and \(N=24\), define
\[
 L_N(x)=\sum_{j=0}^{N}\frac{x^{2j}}{(2j)!},\qquad
 t_{N+1}(x)=\frac{x^{2N+2}}{(2N+2)!},
\]
and
\[
 q_N(x)=\frac{x^2}{(2N+3)(2N+4)}.
\]
The positive Taylor series and the decreasing ratio of its omitted terms
give the rational enclosure
\begin{equation}
 L_N(x)\leq\cosh x\leq
 L_N(x)+\frac{t_{N+1}(x)}{1-q_N(x)}.
 \label{coshinterval}
\end{equation}
All arguments appearing in \(Y\) and \(Z\) are rational, so interval
arithmetic using \eqref{coshinterval}, followed by
\eqref{traceformula}, is exact over \(\mathbb Q\).  It yields the
following strictly positive margins:

\begin{center}
\begin{tabular}{@{}ll@{}}
\toprule
Certified comparison & Rational-interval margin \\
\midrule
\(\trnorm{Y}^2>(1.178)^2\) & greater than \(2.72327\times10^{-4}\) \\
\((1.1402)^2>\trnorm{T}^2\) & exactly \(5.604\times10^{-5}\) \\
\((1.2153)^2>\trnorm{Z}^2\) & greater than \(3.18739\times10^{-4}\) \\
\bottomrule
\end{tabular}
\end{center}

The determinant intervals are positive as well:
\(\det Y>0.08668\) and \(\det Z>0.03910\), justifying removal of the
absolute values in \eqref{traceformula}.  The included
\texttt{verify\_counterexample.py} performs every enclosure and comparison
with Python's exact \texttt{Fraction} arithmetic; floating point is used
only to print readable decimal summaries.

For orientation, ordinary numerical evaluation gives
\[
 \trnorm{Y}\approx1.1781155833,\quad
 \trnorm{T}\approx1.1401754251,\quad
 \trnorm{Z}\approx1.2151688567,
\]
and an actual gap of approximately \(0.0011933767\).

\section{Why the counterexample works}

After conjugating away the geometric term, the conjecture asks a
quadratic interpolation inequality for the Schur multipliers
\(T_{ij}\mapsto\cosh(s(h_i-k_j))T_{ij}\).  Scalar Heinz--Heron
interpolation controls each entry separately, but a unitarily invariant
norm also couples the entries through singular values.  In dimension two,
formula \eqref{traceformula} makes the extra coupling explicit: the
determinant contribution bends differently at the intermediate parameter
\(s=0.7\) than at \(s=0\) and \(s=1\).  Independent left and right spectra
create four distinct logarithmic gaps, and the intermediate trace norm
lies above the conjectured quadratic chord.

This also explains why the Hilbert--Schmidt evidence motivating the
conjecture does not settle it: the Frobenius part is entrywise quadratic,
whereas the trace norm contains the additional determinant term.

\section{Provenance, novelty, and scope}

The exact 2013 source PDF, the 2017 target PDF, and a related 2018 paper
are included.  Searches by arXiv id, title, exact conjecture wording,
formula fragments, author/title citations, and the hyperbolic-kernel
formulation found no prior counterexample or explicit resolution of
Zou's Conjecture (4.2).  This is a bounded novelty check, not a claim that
the entire literature has been exhaustively certified.

The result has deliberately narrow scope:
\begin{itemize}[leftmargin=*]
\item it refutes the universal unitarily invariant norm conjecture using
the trace norm already in dimension two;
\item it does not contradict Zou's Hilbert--Schmidt theorem;
\item it does not concern the separate positive-definite-function question,
which arXiv:1709.09813 answers negatively;
\item it corrects, but does not rely on, the crossed-exponent transcription
in arXiv:1709.09813.
\end{itemize}

\section{References}

\begin{enumerate}[leftmargin=*]
\item L.~Zou, \emph{Inequalities related to Heinz and Heron means},
J. Math. Inequal. 7 (2013), 389--397, doi:10.7153/jmi-07-34.
\item Y.~Kapil, C.~Conde, M.~S.~Moslehian, M.~Singh, and M.~Sababheh,
\emph{Norm inequalities related to the Heron and Heinz means},
Mediterr. J. Math. 14 (2017), Article 213; arXiv:1709.09813.
\item F.~Gao and X.~Ma, \emph{Norm inequalities related to the Heinz
means}, J. Inequal. Appl. 2018 (2018), Article 303.
\end{enumerate}

\end{document}
